\section{Research Background and Global Challenge: From Aesthetics to Emotion–Behavior Alignment}

The release of ChatGPT in November 2022 garnered widespread global attention, showcasing the immense potential of Generative Artificial Intelligence (GenAI) to the world. According to forecasts by IDC and Morgan Stanley, the global GenAI market is projected to reach $145.4 billion by 2027, with related revenues potentially reaching approximately $1.1 trillion by 2028~\citep{JiShuGeXinYinLingWeiLaiShengChengShiAISuZaoHeXinFaZhanYinQing2024,GenAIRevenueCould2025}. GenAI is driving paradigm shifts across various industries. Core technologies such as diffusion models, convolutional neural networks, and generative adversarial networks have enabled the automated generation of text, images, and even videos, presenting unprecedented opportunities for innovation and efficiency \citep{caoComprehensiveSurveyAIgenerated2023,wuAIGCempoweredMethodologyProduct2024}.

The creative design industry is actively embracing this wave. A key application of GenAI, AI-Generated Content (AIGC)—defined as content generated by AI directly from natural language instructions, reference images, or structural templates—has injected new vitality into creative practices by supporting rapid iteration, creative exploration, and design optimization \citep{caoComprehensiveSurveyAIgenerated2023,liUnveilingComplexityDesigners2025,wuAIGCempoweredMethodologyProduct2024}. However, GenAI-driven design still faces several critical challenges. Existing research has paid limited attention to the challenges faced by non-expert users and has inadequately addressed the integration of key design principles such as User-Centered Design (UCD) \citep{maoFactorsInfluencingUser2023,xuEveryoneArtistStudy2023}. In this context, emotional alignment has emerged as a critical issue to bridge the performance gap in GenAI-driven product design\citep{bukerFosteringSelfefficacyUsability2024,demirbilekProductDesignSemantics2003}.

In the field of computer science, advancements in Affective Computing theory and Natural Language Processing (NLP) models have laid the technological foundation for the recognition and expression of emotional information \citep{devlinBERTPretrainingDeep2019,kleine-cosackRecognitionSimulationEmotions2006,liuSentimentAnalysisSubjectivity2010,liuRoBERTaRobustlyOptimized2019,picardAffectiveComputing1997,vaswaniAttentionAllYou2017}. However, how to effectively manifest these capabilities in specific design contexts remains an open question. Therefore, to further unlock the immense potential of GenAI design and enable its application in broader design fields such as product design and customized consumer goods, researchers must ensure semantic correctness while achieving effective alignment between generated outputs and users' stylistic preferences and emotional expectations. This is particularly crucial in multimodal and interactive design scenarios \citep{liWhereAreWe2024,shalev-arkushinImageRAGDynamicImage2025}.

Colour, as a key component of emotional expression in visual information, offers higher interpretability and stronger manipulability in inducing visual emotions compared to more abstract dimensions like style or semantics \citep{kobayashiAimMethodColor1981,ouCrossculturalComparisonColour2012,valdezEffectsColorEmotions1994}. In recent years, research in fields such as model training and AIGC generation strategy optimization has noted the significant role of colour variables. For instance, colour features are often considered key indicators in tasks like image segmentation, aesthetic assessment, and emotion recognition \citep{nieDeepLearningbasedMachine2023,yuArtificialIntelligenceAIdriven2025}. Studies in cognitive psychology have also begun to investigate the colour-emotion associations in text-to-image models. This body of research indicates that colour can not only reliably evoke emotional responses but also that its parametric features (e.g., hue, brightness, saturation) provide operable pathways for manipulation in design practice. Nevertheless, despite these advances, existing explorations remain fragmented across domains and have yet to converge into a generalizable mechanism that systematically links target emotions and usage contexts to controllable colour variables in generative design.

As shown in Figure 1, this study focuses on establishing a controllable and testable pipeline that links the features of GenAI-designed objects, the alignment of human-AI affective perception, design attractiveness, and user adoption intention, validated against a human baseline.

\begin{figure}[htbp]

  \centering

  % TODO: insert figure file

  % \includegraphics[width=0.9\textwidth]{figures/your_figure_file}

  \caption{Figure 1: Background and Missing PieceThe Research Gap: Measurement Instability and Generative Bias in Emotion Control.}

  \label{fig:figure_1_background_and_missing_piecethe}

\end{figure}

\section{Research Gap: Measurement Instability and Generative Bias in Emotion Control}

To address the "missing piece" identified in Section 1.1—namely, the disconnect between the outcomes produced by GenAI-driven applications and their genuine adoption by users—a substantial body of recent research has concentrated on enhancing the quality of model-generated results. These efforts have primarily focused on the refinement of model training processes and the development of strategies for AI image enhancement. In the domain of image generation, Re-Imagen \citep{chenReimagenRetrievalaugmentedTexttoimage2022} and RDMs \citep{blattmannRetrievalaugmentedDiffusionModels2022} have demonstrated the effectiveness of retrieval-augmented generation \citep{blattmannRetrievalaugmentedDiffusionModels2022,chenReimagenRetrievalaugmentedTexttoimage2022}. More recently, RealRAG \citep{lyuRealRAGRetrievalaugmentedRealistic2025} has further improved the expressiveness of text-to-image models for rare and fine-grained entities \citep{lyuRealRAGRetrievalaugmentedRealistic2025}. These studies have made notable contributions to improving the quality and consistency of generated outputs, but such advances alone are insufficient for GenAI applications—such as text-to-image models—to be truly adopted in design practice. What is further required is the ability to interpret the implicit affective expectations embedded in user queries and to reflect them in the generated results, thereby enabling outputs that resonate emotionally with consumers and support innovative, personalized product design \citep{yangEmotionallyIntelligentFashion2021}.

While these challenges surface most clearly in design practice, they can be further traced to specific technical layers—most notably in natural language processing and visual generation,. Existing studies indicate that models capture dominant emotions but struggle with complex, mixed, or context-dependent affective states \citep{acheampongTransformerModelsTextbased2021,finetLimitsAIUnderstanding2025}. In visual generation, AI tends to output positive emotional expressions, exhibiting a significant bias in the representation of negative or ambiguous emotions \citep{carrascoLevelAgreementEmotions2024,zhangDecodingEmotionalResponses2024}. Furthermore, the coverage and balance of emotional data resources are limited; for instance, the ArtEmis v2.0 dataset reveals a highly imbalanced distribution of emotion categories \citep{mohamedItOkayNot2022}, which further weakens the model's training effectiveness fatoresin the emotional dimension \citep{mohamedItOkayNot2022}. Consequently, even though retrieval-augmented generation has made significant strides in semantic accuracy and content diversity, a distinct gap remains between understanding the user's emotional intent and achieving fine-grained emotional control over the generated results.

Against this backdrop, some studies have begun to explore new technical approaches to enhance the emotional alignment and user adaptivity of generated outcomes. For example, Yu, Cheng, and Tian (2025) proposed a multimodal AI art generation strategy based on stable diffusion and BERT, aiming to improve controllability, emotional alignment, and user adaptivity in AI art synthesis \citep{yuArtificialIntelligenceAIdriven2025}. Hou et al. (2025) explores more detailed manipulation of semantic and emotional features within the generation process by introducing conditional controls into diffusion models \citep{houGenColorGenerativeColorconcept2025}. Notably, these recent works have all integrated color-emotion mapping into AI generation strategies as a crucial method for enhancing the quality of the results. These recent works have, in different ways, converged on integrating color–emotion mapping into AI generation strategies as a crucial pathway for improving the performance of generated results. The explanatory power and operability of color in evoking visual emotions have been substantiated in concrete GenAI production practices, as surveyed in Sections 1.1 and2.1. However, these discussions and limitations suggest that most approaches rely on specific datasets or experimental environments and have yet to propose stable emotional variables that can be operationalized in design practice. Therefore, building upon existing explorations, this study aims to further focus on the design context of consumer product customization, examining the color-emotion mapping mechanism as a controllable variable.

\section{Research Object: Floral Gifting as an Emotion-Oriented Design Scenario}

In design contexts, a user's emotional resonance directly influences their intention to adopt a design solution—a concept originally derived from research in the consumer products domain. Color plays a pivotal role in shaping individuals' emotional perception and cognitive appraisal of objects \citep{elliotColorPsychologyEffects2014,yaoVisualPerformancePainting2022}, and this emotional response, elicited by the design, is a direct determinant of whether the final product can fulfill the deeper-level needs of the consumer. Consequently, in numerous consumption behaviors contingent upon emotional resonance, the congruence among visual design, intended emotion, and the user's ultimate choice is of decisive importance.

To systematically address the aforementioned challenges, this study selects floral gift design as a representative test scenario. As a quintessential artistic consumer product, a floral gift is distinguished by its dual possession of aesthetic value and emotional functionality. Unlike traditional floristry or landscape horticulture, modern floral gifts emphasize "miniaturization, personalization, and contextual fit," encompassing a variety of forms such as fresh flower bouquets, preserved flower gift boxes, and other gift-oriented floral compositions. These products are extensively utilized in scenarios such as birthdays, anniversaries, festivals, and business correspondence, catering to a diverse user base and thus possessing significant universality and research value.

The design process of a floral gift is highly dependent on emotional expression. Its ultimate emotional effect is a composite of multiple factors: first, the symbolic meaning carried by the flowers themselves, i.e., "floriography," has been used since antiquity to convey interpersonal sentiments and cultural symbols \citep{dietzCompleteLanguageFlowers2022,JiaJunZhiWuYiXiangYanJiu2011}; second, the psychological effects induced by color combinations. Extensive research has demonstrated that different color palettes can evoke specific emotions and exert a significant influence on cognition and decision-making. The interplay of these two factors renders a floral gift not merely a materialized commodity but a medium that carries emotional information and social significance.

Based on these considerations, this study selects floral gifts as its research object: on the one hand, it represents a quintessential scenario within consumer product design that is most heavily reliant on emotional resonance; on the other hand, its high sensitivity to color-emotion alignment makes it an ideal case for validating whether generative design can achieve emotional alignment. This provides a clear context and experimental foundation for the subsequent research objectives and research questions.

\section{Research Study Overview and Hypotheses}

This study aims to construct an integrative framework that connects color features in design, user-perceived emotional fit, and the final behavioral outcome (i.e., adoption intention). Extant literature and practice indicate that the success of a design depends not merely on its aesthetic appearance, but more critically on whether its visual elements can accurately convey and evoke the intended emotions. When a design achieves consistency with user needs on an emotional level, its perceived appeal and the user's adoption intention are significantly enhanced. Therefore, the objective of this study is to explore a controllable generation method for context-sensitive visual design of consumer products, validated by user behavior, to address the shortcomings of current generative AI in achieving emotional alignment. Accordingly, this study is expected to make the following contributions:

\begin{itemize}

  \item Theoretical contribution: Develop an integrative framework that connects color features, emotional fit, design appeal, and adoption intention, enriching the understanding of the “emotion–aesthetics–behavior” pathway in consumer design contexts.

  \item Empirical contribution: Establish a human baseline for color–emotion mapping and provide comparative evidence on the biases and deficiencies of AIGC-generated designs.

  \item Methodological contribution: Propose a user-centered approach to evaluating emotional alignment in design outputs, linking perceptual and behavioral measures.

  \item Practical contribution: Offer actionable suggestions for researchers and system developers, particularly regarding the incorporation of emotion-informed asset libraries and generation strategies into retrieval–generation pipelines for GenAI-driven customized design.

\end{itemize}

Building on the research gaps identified above, this thesis addresses a set of research questions that move from a fundamental investigation of colour–emotion relations to their application and validation in GenAI-supported design. As shown in Figure 2, the research is organised into three progressive levels:

\begin{itemize}

  \item RQ1: In a contextualised floral gift design scenario, how do colour-related features influence users’ perceived emotional fit?

  \item (To address the fundamental question of whether colour can serve as a stable, controllable carrier of intended emotion in consumer-oriented design. It focuses on identifying which colour attributes and palette structures are associated with higher emotional alignment under different gifting contexts.)

  \item RQ1： 在具体语境化的花礼设计情境中，颜色相关特征如何影响用户感知到的情绪契合？

  \item RQ2: How does perceived emotional fit relate to downstream user responses, particularly design appeal and adoption intention?

  \item (To examine the mechanism through which emotion alignment shapes behavioural outcomes. Specifically, it asks whether emotional fit directly contributes to adoption intention and whether this relationship is mediated by users’ aesthetic evaluation of the design.)

  \item RQ2： 感知到的情绪契合如何与后续的用户反应相关，尤其是设计吸引力与采纳意愿？

  \item RQ3: To what extent can these colour–emotion relationships be applied to the evaluation and improvement of GenAI-generated design outputs, in comparison with human-created designs?

  \item (To extend the fundamental and mechanistic findings into an application-oriented setting. It investigates whether the established human baseline can reveal systematic biases or deficiencies in AIGC-generated designs, and whether such evidence can support more controllable and emotionally aligned generative design strategies.)

  \item RQ3： 与人类创作设计相比，这些色彩—情绪关系在多大程度上可以被应用于 GenAI 生成设计的评估与改进？

\end{itemize}

\begin{figure}[htbp]

  \centering

  % TODO: insert figure file

  % \includegraphics[width=0.9\textwidth]{figures/your_figure_file}

  \caption{Figure 2: Research QuestionsThe Upstream and Downstream Pathways of Color–Emotion Alignment in Generative Design}

  \label{fig:figure_2_research_questionsthe_upstream}

\end{figure}

Three research questions establish a coherent progression from foundational inquiry to practical validation. RQ1 identifies the upstream relationship between colour and emotional alignment; RQ2 examines the downstream pathway from emotional fit to aesthetic and behavioural responses; and RQ3 applies these findings to the evaluation of GenAI-generated design, linking basic research to real design implementation.

Figure 2 illustrates the progressive relationship among these three research questions. Together, they address the missing piece between generative design research and its application in actual consumer-oriented design practice.

To address the first two research questions, this study further proposes two research hypotheses:

H1: In a contextualised floral gift design scenario, colour-related features are significantly associated with users’ perceived emotional alignment.

H1： 在具体语境化的花礼设计情境中，颜色相关特征与用户感知到的情绪契合显著相关。

H2: Perceived emotional alignment is positively associated with design appeal and adoption intention, and the relationship between perceived emotional alignment and adoption intention is partially mediated by design appeal.

H2： 感知到的情绪契合与设计吸引力及采纳意愿呈正相关，且情绪契合对采纳意愿的影响部分通过设计吸引力实现。

Figure 3 shows how these two hypotheses support the operational structure of the present study. Specifically, it translates the relationships outlined in RQ1 and RQ2 into a testable analytical framework linking colour-related features, emotional fit, design appeal, and adoption intention. In contrast, RQ3 is addressed through the comparative evaluation of human-created and GenAI-generated design outputs, rather than through a separate formal hypothesis.

\begin{figure}[htbp]

  \centering

  % TODO: insert figure file

  % \includegraphics[width=0.9\textwidth]{figures/your_figure_file}

  \caption{Figure 3: Research HypothesesLinking Color-Induced Emotional Fit and Adoption Mechanisms toward Feasible Design Outputs}

  \label{fig:figure_3_research_hypotheseslinking_colo}

\end{figure}
